% Complete proofs for ordinal collective choice and binary majority.

\begin{annexitem}{annex:p17.social-majority}{thm:p17-social-majority}{Pair counts, the cycle and the binary case}
\begin{proof}
For distinct \(x,y\), every strict input ordering contributes to exactly
one of \(V_P(x,y)\) and \(V_P(y,x)\), so their sum is \(n\).
The two integer counts are comparable. The weak majority relation is
therefore complete, and its strict part is precisely strict numerical
inequality between the counts. The explicitly included diagonal makes it
reflexive. If all voters rank \(x\) above \(y\), the counts are \(n,0\);
since \(n\geq1\), the social comparison is strict. The definition uses only
the individual comparisons of the one pair.

In the displayed profile, voters 1 and 3 prefer \(a\) to \(b\),
voters 1 and 2 prefer \(b\) to \(c\), and voters 2 and 3 prefer \(c\) to
\(a\). The remaining voter reverses each comparison. Hence each asserted
strict majority has count two against one. Transitivity would imply
\(a\succ^M_Pc\) from the first two comparisons, contrary to the third.

On a two-element set, every complete reflexive relation is transitive.
Indeed, if \(x\succeq y\succeq z\), then either \(x=z\), when reflexivity
gives the conclusion, or \(x\ne z\). In the latter case \(y\) equals
one of \(x,z\), and one of the two premises is already \(x\succeq z\).
Thus majority is a social welfare function in this domain, including
equal-count ties when the electorate has even size.

For three voters, fix any proposed dictator \(j\). Let that voter rank
\(a\) above \(b\) and the two others rank \(b\) above \(a\).
The social strict comparison is opposite to \(j\)'s comparison.
This constructs a refuting profile for each of the three possible fixed
dictators, proving nondictatorship.
\end{proof}
\end{annexitem}

\begin{annexitem}{annex:p17.social-aggregation-rules}{thm:p17-social-aggregation-rules}{Defined rankings and the changed victory score}
\begin{proof}
A constant rule returns a weak preference by its definition.
Any real-valued score \(u\) defines a complete transitive relation by
\(x\succeq y\) exactly when \(u(x)\geq u(y)\): completeness and
transitivity follow from those of the real order, including ties.
This proves the assertion for victory scores and rank scores.
The threshold procedure returns one of the two strict binary orderings,
whose weak closures are complete and transitive.

Reversing a strict ordering gives another strict ordering. Irreflexivity
and comparison of distinct pairs are immediate. If its reversal puts
\(x\) above \(y\) and \(y\) above \(z\), then the original puts
\(z\) above \(y\) above \(x\); original transitivity puts \(z\) above
\(x\), as required. Adding the diagonal gives a weak preference.

For the minimum-disagreement procedure, a finite nonempty \(X\) has
finitely many strict orderings and at least one. For each ordering,
the displayed set of triples is finite, so its cardinality is a
nonnegative integer. A minimum over this finite nonempty candidate set
exists. The fixed priority of candidates selects exactly one minimizer.
The most-common-ranking procedure likewise maximizes a finite integer
count. Because \(n\geq1\), its maximum is positive, so a maximizing
ordering occurs among the inputs. The fixed priority resolves every tie.
Both procedures therefore return the weak closure of a strict ordering.

For the Pareto assertion about rank scores, suppose every voter ranks
\(x\) above \(y\). Then \(r_i(x)<r_i(y)\) for every \(i\), so
\(w_{r_i(x)}>w_{r_i(y)}\). Summing the finitely many strict positive
differences gives \(B_P(x)>B_P(y)\).

For victory scores under the same unanimity premise, \(x\) defeats \(y\).
If \(y\) defeats a third alternative \(z\), every voter preferring \(y\)
to \(z\) also prefers \(x\) to \(z\), by individual transitivity.
Consequently \(x\) also defeats \(z\). Thus the set of victories of
\(y\) is contained in the victories of \(x\), and \(y\) itself is an
additional victory of \(x\). The former set contains neither \(x\) nor
\(y\). Hence \(C_P(x)\geq C_P(y)+1\), proving the strict social comparison.

At the cyclic profile each alternative defeats one other alternative and
loses to the remaining one, giving all scores one. At \(P'\), the
individual orderings are \(a>b>c\), \(b>a>c\), and \(c>a>b\).
The votes favor \(a\) over \(b\) by voters 1 and 3; \(a\) over \(c\)
by voters 1 and 2; and \(b\) over \(c\) by voters 1 and 2.
The victory scores are therefore \(2,1,0\) as asserted.
Voter 2's change preserves \(b>a\), and neither other voter changes.
Initially \(b\succeq_Pa\) by a tied score; finally that weak comparison
is false. Thus identical individual comparisons of this pair need not
give identical social comparisons under the victory-score procedure.
\end{proof}
\end{annexitem}

\begin{annexitem}{annex:p17.social-extreme-alternative}{thm:p17-social-extreme-alternative}{The extreme-alternative argument with possible social ties}
\begin{proof}
Suppose \(b\) is not strictly above every other alternative and is not
strictly below every other alternative in the complete social relation.
Then the sets
\[
 A=\{x\in X\setminus\{b\}:x\succeq_Pb\},\qquad
 C=\{x\in X\setminus\{b\}:b\succeq_Px\}
\]
are both nonempty. To justify the first assertion, failure of
\(b\succ_Px\), together with completeness, implies \(x\succeq_Pb\);
the other assertion is the reversed argument.

There exist distinct \(a\in A\) and \(c\in C\). Otherwise nonemptiness
and the absence of any distinct pair would force \(A=C=\{d\}\) for
one alternative \(d\). Since \(X\) contains at least three members,
there is \(e\notin\{b,d\}\). Completeness puts \(e\) in \(A\) or \(C\),
contradicting those singleton identities. Thus
\[
                         a\succeq_Pb\succeq_Pc
\]
holds for distinct \(a,b,c\), even if one or both comparisons are ties.

Modify each individual ordering to place \(c\) above \(a\) while retaining
whether \(b\) is at its top or bottom. For example, among alternatives
other than \(b\), remove \(c\) and insert it immediately before \(a\).
All other comparisons in that remaining order may be retained.
This is a strict ordering, including for an infinite alternative set:
elements above \(a\) remain above the inserted \(c\), and \(c\) precedes
\(a\) and all its successors. With \(b\) left at its old extreme, neither
the individual \(a,b\) comparison nor the \(b,c\) comparison changes.
Unrestricted domain makes the modified profile \(P'\) admissible.

IIA preserves the social weak comparisons
\(a\succeq_{P'}b\succeq_{P'}c\). Transitivity gives
\(a\succeq_{P'}c\). But every voter now strictly prefers \(c\) to \(a\);
Pareto gives \(c\succ_{P'}a\), which by definition forbids
\(a\succeq_{P'}c\). This contradiction proves that \(b\) is strictly
top or strictly bottom socially. Both alternatives cannot hold when
there is another alternative. Either conclusion also excludes any tie
between \(b\) and another alternative.
\end{proof}
\end{annexitem}

\begin{annexitem}{annex:p17.social-arrow}{thm:p17-social-arrow}{A pivotal voter and all pair comparisons}
\begin{proof}
If the profile domain is empty, every statement quantified over profiles
is vacuous, and any of the \(n\geq1\) voters satisfies the definition of
dictator. Otherwise choose a profile, fix \(b\in X\), and retain its
individual orderings on \(X\setminus\{b\}\) as reference orderings.
In the ordinary finite-alternative case this domain is nonempty because
a finite set can be ordered. The remaining argument applies to any
nonempty strict-profile domain, without requiring \(X\) to be finite.

\paragraph{A finite sequence of extreme profiles.}
For each \(m=0,\ldots,n\), insert \(b\) at the top for voters
\(i\leq m\) and at the bottom for voters \(i>m\), obtaining \(P^m\).
These are strict orderings on all of \(X\). The preceding lemma makes
\(b\) socially strictly top or strictly bottom at each step.
At \(m=0\), unanimity puts every other alternative strictly above \(b\).
At \(m=n\), unanimity puts \(b\) strictly above every other alternative.
Let \(m^*\) be the smallest index with the latter property. It exists
in \(\{1,\ldots,n\}\). Minimality and the extreme-alternative lemma put
\(b\) strictly bottom at \(P^{m^*-1}\).

No monotonicity axiom about these intermediate social positions has been
used: the first top step exists even without assuming that top status
persists at all later steps.

For each \(m\) and each \(x\ne b\), the vector of individual comparisons
of \(b,x\) is determined solely by whether \(i\leq m\).
Changing the reference orderings of pairs not involving \(b\) preserves
that vector. Applying IIA to both orientations of the pair preserves
the strict social comparison. Hence the top/bottom status at each \(m\),
and therefore \(m^*\), is independent of the reference orderings.

\paragraph{Pairs not involving the reference alternative.}
Let \(a,c\) be distinct and different from \(b\). Take any profile \(P\)
at which \(a\succ_{m^*}c\). Form \(P'\) by removing \(b\) from each
ordering, leaving all other comparisons fixed, and inserting it as follows:
at the top for \(i<m^*\), between \(a\) and \(c\) for \(i=m^*\),
and at the bottom for \(i>m^*\).
For the middle voter, insertion immediately before \(c\) places it below
\(a\) and above \(c\); this also defines a strict ordering on infinite \(X\).

The pair \(b,a\) at \(P'\) has exactly the same individual comparisons
as at \(P^{m^*-1}\): the voters before \(m^*\) put \(b\) above \(a\),
and all other voters put \(a\) above \(b\). IIA therefore gives
\(a\succ_{P'}b\).
The pair \(b,c\) has exactly the same individual comparisons as at
\(P^{m^*}\): voters through \(m^*\) put \(b\) above \(c\), and the
others put \(c\) above \(b\). IIA gives \(b\succ_{P'}c\).

The strict part of a complete transitive weak relation is transitive.
Indeed, from \(u\succ v\succ w\), weak transitivity gives \(u\succeq w\);
if \(w\succeq u\), transitivity with \(u\succeq v\) would contradict
\(v\succ w\). Thus \(u\succ w\).
Applying this fact gives \(a\succ_{P'}c\). No individual \(a,c\)
comparison changed from \(P\) to \(P'\), so IIA in both orientations
gives \(a\succ_Pc\). We have proved that voter \(m^*\) determines
every strict comparison of two alternatives other than \(b\).

\paragraph{Pairs involving the reference alternative: the first direction.}
Fix \(a\ne b\) and a profile \(P\) with \(a\succ_{m^*}b\).
Choose \(c\notin\{a,b\}\), possible because \(X\) has at least three
members. Make \(P'\) by moving \(c\) to the top for every voter other
than \(m^*\), and inserting it between \(a,b\) for voter \(m^*\).
Retain all comparisons not involving \(c\), in particular all individual
\(a,b\) comparisons.

For the pivot voter, \(a\succ_{m^*}'c\); these two alternatives differ
from \(b\), so the previous paragraph gives \(a\succ_{P'}c\).
Every voter prefers \(c\) to \(b\), so Pareto gives \(c\succ_{P'}b\).
Strict transitivity and IIA give first \(a\succ_{P'}b\), then
\(a\succ_Pb\).

\paragraph{Pairs involving the reference alternative: the reverse direction.}
Now take a profile \(P\) with \(b\succ_{m^*}a\) and the same choice
of a third alternative \(c\). Move \(c\) to the bottom for every voter
other than \(m^*\), and between \(b,a\) for voter \(m^*\), retaining
all other comparisons. Every voter now prefers \(b\) to \(c\), so
Pareto gives \(b\succ_{P'}c\). The pivot voter prefers \(c\) to \(a\),
so the already proved result for pairs not involving \(b\) gives
\(c\succ_{P'}a\). Strict transitivity and IIA give \(b\succ_Pa\).
This supplies explicitly the reversed case in the source's last step.

\paragraph{One fixed voter determines the entire relation.}
Every distinct pair either avoids \(b\) or includes it, so the preceding
arguments cover every pair and every profile. For each such pair the
pivot's strict ordering picks one direction, and the social strict
comparison agrees. The social weak comparison in the reverse direction
is therefore false. On the diagonal completeness gives reflexivity.
Consequently the complete social relation is precisely the reflexive
closure of voter \(m^*\)'s ordering, at every profile. This is the
required fixed dictator.
\end{proof}
\end{annexitem}

\begin{annexitem}{annex:p17.social-arrow-boundaries}{thm:p17-social-arrow-boundaries}{Four explicit failures of a removed premise}
\begin{proof}
For the first example one may take three alternatives and any positive
finite number of voters. The reversal rule returns a weak preference,
as proved above. Its comparison of each pair depends only on the
designated voter's comparison of that pair, hence it satisfies IIA.
At a unanimous profile it reverses every distinct pair, so violates
Pareto. The same unanimous profile also shows it cannot be dictatorial
in the stated sense: its output disagrees with each possible voter.

For the rank-score example, the preceding theorem proves transitivity,
completeness and Pareto. At \(P\), voter 1 contributes scores
\((2,1,0)\) to \((a,b,c)\), and voter 2 contributes \((0,2,1)\);
the sum is \((2,3,1)\). At \(P'\), the contributions are
\((2,0,1)\) and \((1,2,0)\), summing to \((3,2,1)\).
Voter 1 retains \(a>b\), and voter 2 retains \(b>a\), but the social
strict comparison reverses, so IIA fails. The rule is not dictatorial:
at \(P\), the output \(b>a>c\) differs from voter 1 on \(a,b\) and
from voter 2 on \(a,c\). Thus neither of the two fixed voters can
determine all profiles.

The three-voter majority example is complete, Pareto and pairwise
independent by Theorem~\ref{thm:p17-social-majority}. Its cycle is an
explicit violation of transitivity. It is nondictatorial because, for
each proposed voter, the other two can unanimously oppose that voter's
comparison on a chosen pair; complete strict rankings realizing those
votes can be supplied on the three alternatives.

Finally the same theorem proves that binary majority with three voters
is a complete transitive social relation, follows unanimity, is pairwise
independent, and is nondictatorial. Its alternative set has only two
members, locating precisely the premise that changes.
\end{proof}
\end{annexitem}

\begin{annexitem}{annex:p17.social-may}{thm:p17-social-may}{Balanced profiles determine the unique binary rule}
\begin{proof}
First suppose \(f\) has the three properties. At a profile \(s\) with
equal numbers of coordinates \(1\) and \(-1\), there is a permutation
of voters sending \(s\) to \(-s\): pair the positive and negative
positions and interchange each pair, leaving the zero positions fixed.
Anonymity gives \(f(s)=f(-s)\), while neutrality gives
\(f(-s)=-f(s)\). The only member of \(\{-1,0,1\}\) equal to its
negative is zero. Hence every such balanced profile has output zero.
This includes the all-indifferent profile.

Suppose next that \(s\) has more positive than negative coordinates.
Replace enough of its positive coordinates by zero to make their number
equal to its number of negative coordinates, obtaining \(t\).
The changed set is finite and nonempty. The balanced case gives
\(f(t)=0\). Restore the changed coordinates one at a time from zero
to one. The first restoration has a nonnegative previous output, so
positive responsiveness gives output one. Each later restoration again
starts from a nonnegative output and therefore keeps output one.
After all restorations we obtain \(f(s)=1\).
If \(s\) has more negative than positive coordinates, the result just
proved applied to \(-s\), followed by neutrality, gives \(f(s)=-1\).
These three cases exhaust all profiles and yield the asserted formula.

Conversely, define \(f(s)\) by that formula. The finite sum is unchanged
by any permutation, so the rule is anonymous. Negation reverses the
sum's sign, giving neutrality. Increasing one coordinate by one
increases the integer sum by exactly one. If its former sign was
nonnegative, its new value is strictly positive, proving the positive
direction of responsiveness. Decreasing one coordinate by one from
a nonpositive sum makes it strictly negative, proving the symmetric
direction as well. All outputs encode complete transitive preferences
on the two alternatives, including the zero-output tie.
This proves existence and uniqueness on the stated weak-input domain.
\end{proof}
\end{annexitem}

\begin{annexitem}{annex:p17.social-may-boundaries}{thm:p17-social-may-boundaries}{Separating the binary axioms}
\begin{proof}
For the weighted rule, negating the input negates \(2s_1+s_2\), so
neutrality holds. Raising one coordinate by one raises this weighted
sum by the corresponding positive weight, either two or one. From
a nonnegative prior weighted sum the result is strictly positive.
Similarly a coordinate decrease from a nonpositive sum makes it
strictly negative. Thus both directions of positive responsiveness hold.
But \(f(1,-1)=1\), whereas after interchanging the voters
\(f(-1,1)=-1\). Anonymity fails.

The constant-one rule ignores voter identities, so is anonymous.
Every upward change has output one, as the positive responsiveness
conclusion requires. The symmetric downward condition has a false
premise, because the original output is never nonpositive.
Hence the full responsiveness condition holds. Neutrality fails since
both \(f(s)\) and \(f(-s)\) equal one rather than being negatives.

The constant-zero rule is anonymous and neutral. At \(s=(0,0)\),
raise the first coordinate to one. The original output is nonnegative,
but the new output is still zero, violating positive responsiveness.
Each example therefore keeps exactly the two claimed conditions while
refuting the third.
\end{proof}
\end{annexitem}
